% !TEX TS-program = pdflatex
% !TeX program = pdflatex
% !TEX encoding = UTF-8
% !TEX spellcheck = en_US

\documentclass[xcolor=table]{beamer}

\usepackage{../extra/beamer/karimnlp}


%\usepackage{fullpage}
%\usepackage[left=2.8cm,right=2.2cm,top=2 cm,bottom=2 cm]{geometry}
\setbeamersize{text margin left=10pt,text margin right=10pt}



%\usepackage[T1]{fontenc}
\usepackage{textcomp}
\usepackage[utf8]{inputenc}
%\usepackage[french,english]{babel}
\usepackage{arabtex}
\usepackage{txfonts}
\usepackage{tipa}
\usepackage[]{graphicx}
\usepackage{multirow}
\usepackage{hyperref}
\usepackage{colortbl}
\usepackage{tabularray}
\usepackage{listingsutf8}
\usepackage{wrapfig}
\usepackage{multicol}
\usepackage[export]{adjustbox} %for images in table, also for frame
\usepackage[many]{tcolorbox}
\usepackage{wasysym}
\usepackage[lined]{algorithm2e}
\usepackage{alltt} %verbatim with commands
\usepackage{longtable}
\usepackage{tabu}

\usepackage{CJKutf8}

\usepackage{amsmath,amssymb, amsfonts} 

%\usepackage{newtxtext,newtxmath}


%\usepackage{acolor}


\definecolor{lightblue}{HTML}{D0D2FF}
\definecolor{lightyellow}{HTML}{FFFFAA}
\definecolor{darkblue}{HTML}{0000BB}
\definecolor{olivegreen}{HTML}{006600}
\definecolor{darkgreen}{HTML}{008B45} %009B55
\definecolor{violet}{HTML}{6600CC}
\definecolor{deeppink}{HTML}{FF1493}
\definecolor{orangey}{HTML}{FFBB00}


\makeatletter
\def\beamer@calltheme#1#2#3{%
	\def\beamer@themelist{#2}
	\@for\beamer@themename:=\beamer@themelist\do
	{\IfFileExists{\beamer@themelocation/#3\beamer@themename.sty}{%
	\usepackage[{#1}]{\beamer@themelocation/#3\beamer@themename}}{\usepackage[{#1}]{#3\beamer@themename}}
		
	}%
}

\def\usefolder#1{
	\def\beamer@themelocation{#1}
}
\def\beamer@themelocation{}
\makeatother

\usefolder{../../extra/beamer-themes}

\usetheme{Karim} % Antibes Boadilla Warsaw Karim

\beamertemplatenavigationsymbolsempty

\let\oldcite\cite
\renewcommand{\cite}[1]{{\bfseries\color{blue}\oldcite{#1}}}

\tcbuselibrary{listings}


%\renewcommand{\baselinestretch}{1.5}

\def\supit#1{\raisebox{0.8ex}{\small\it #1}\hspace{0.05em}}

\AtBeginSection{%
	\begin{frame}
		\sectionpage
	\end{frame}
}

\newcommand{\rottext}[2]{%
	\rotatebox{90}{%
	\begin{minipage}{#1}%
		\raggedleft#2%
	\end{minipage}%
	}%
}


\institute{ %
Laboratoire de la Communication dans les Systèmes Informatiques (LCSI)
\par
École  nationale Supérieure d'Informatique (ESI), Algiers, Algeria
}
\author[Abdelkrime Aries (2023/2024)] %
{Abdelkrime Aries}
%\titlegraphic{\includegraphics[height=1cm]{../img/esi-logo.png}%\hspace*{4.75cm}~


\date{Academic year: 2023/2024} %\today

\titlegraphic{%
	\includegraphics[height=1cm]{../../img/misc/esi-logo.png}%
	\hspace{2cm}%
	\includegraphics[height=1cm]{../../img/misc/lcsi-logo.png}%
	\hspace{2cm}%
	\includegraphics[height=1.5cm]{../../img/misc/esi.nlp-logo.png}%
}

%\setbeamertemplate{headline}{}

\newcommand{\kurl}[1]{{\scriptsize\bfseries\color{orangey}\url{#1}}}

\newcommand{\keyword}[1]{\textcolor{red}{\bfseries\itshape #1}}
\newcommand{\expword}[1]{\textcolor{olivegreen}{#1}}
\newcommand{\optword}[1]{\textcolor{violet}{\bfseries #1}}

\makeatletter
\newcommand\mysphere{%
	\parbox[t]{10pt}{\raisebox{0.2pt}{\beamer@usesphere{item projected}{bigsphere}}}}
\makeatother

%\let\oldtabular\tabular
%\let\endoldtabular\endtabular
%\renewenvironment{tabular}{\rowcolors{2}{white}{lightblue}\oldtabular\rowcolor{blue}}{\endoldtabular}


%\NoAutoSpacing %french autospacing after ":"

\def\graphpath{}

\newcommand{\changegraphpath}[1]{\def\graphpath{#1}}


\newcommand{\vgraphpage}[2][.84\textheight]{%
%	\begin{center}%
		\includegraphics[height=#1]{\graphpath #2}%
%	\end{center}%
}

\newcommand{\hgraphpage}[2][\textwidth]{%
%	\begin{center}%
		\includegraphics[width=#1]{\graphpath #2}%
%	\end{center}%
}

\newcommand{\graphpage}[2][]{%
	\includegraphics[#1]{\graphpath #2}%
}

\bibliographystyle{apalike}

\newcommand{\insertbibliography}[2]{
	\appendix
	\section*{Bibliography}
	\nocite{#2}
%	\makeatletter % to change template
%	\setbeamertemplate{headline}[default] % not mandatory, but I though it was better to set it blank
%	\def\beamer@entrycode{\vspace*{-\headheight}} % here is the part we are interested in :)
%	\makeatother
	\begin{multicols*}{2}[\frametitle{\insertsection} \usebeamertemplate*{frametitle}]%\usebeamertemplate*{frametitle}\frametitle{Références}
		\tiny
		\bibliography{#1}
	\end{multicols*}
}

\definecolor{my-grey}{RGB}{233, 233, 233}

\newcommand{\insertlicence}{
	\begin{frame}[plain]
	\frametitle{License: CC-BY 4.0}

	\begin{tcolorbox}[colback=cyan,
		colframe=cyan,  
		arc=0pt,outer arc=0pt,
		valign=top, 
		halign=center,
		width=\textwidth]
		
		\includegraphics[width=.5cm]{../../img/licence/cc_icon_white_x2.png}
		\includegraphics[width=.5cm]{../../img/licence/attribution_icon_white_x2.png}
		
		\color{white}
		\bfseries Attribution 4.0 International (CC BY 4.0) \\
		\tiny \url{https://creativecommons.org/licenses/by/4.0/deed.en}
		
	\end{tcolorbox}\vspace{-.5cm}
	\begin{tcolorbox}[colback=my-grey,
		colframe=my-grey,  
		center, arc=0pt,outer arc=0pt,
		valign=top, 
		halign=left,
		width=\textwidth]
		
		\tiny
		
		\begin{center}
			\bfseries\Large
			You are free to:
		\end{center}
		
		\begin{minipage}{0.83\textwidth}
			\begin{itemize}
				\item[] \textbf{Share} — copy and redistribute the material in any medium or format
				\item[] \textbf{Adapt} — remix, transform, and build upon the material
				for any purpose, even commercially.
			\end{itemize}
		
			\vspace{12pt}
		
			\textit{The licensor cannot revoke these freedoms as long as you follow the license terms.}
		\end{minipage}
		\begin{minipage}{0.15\textwidth}
			\includegraphics[width=\textwidth]{../../img/licence/FreeCulturalWorks_seal_x2.jpg}
		\end{minipage}
	
		
		\begin{center}
			\bfseries\Large
			Under the following terms:
		\end{center}
		
		\begin{itemize}
			\item[] \textbf{Attribution} — You must give appropriate credit, provide a link to the license, and indicate if changes were made. You may do so in any reasonable manner, but not in any way that suggests the licensor endorses you or your use.
			\item[] \textbf{No additional restrictions} — You may not apply legal terms or technological measures that legally restrict others from doing anything the license permits.
		\end{itemize}
		
	\end{tcolorbox}
	
%	\begin{center}
%		\bfseries Attribution 4.0 International (CC BY 4.0)
%		\url{https://creativecommons.org/licenses/by/4.0/deed.fr}
%	\end{center}

%	\tiny
%
%	Vous êtes autorisé à : 
%	\begin{itemize}
%		\item \textbf{Partager} — copier, distribuer et communiquer le matériel par tous moyens et sous tous formats
%		\item \textbf{Adapter} — remixer, transformer et créer à partir du matériel
%	\end{itemize}
%	
%	Selon les conditions suivantes : 
%	\begin{itemize}
%		\item \textbf{Attribution} — Vous devez créditer l'Œuvre, intégrer un lien vers la licence et indiquer si des modifications ont été effectuées à l'Oeuvre. Vous devez indiquer ces informations par tous les moyens raisonnables, sans toutefois suggérer que l'Offrant vous soutient ou soutient la façon dont vous avez utilisé son Oeuvre.
%		\item \textbf{Pas d'Utilisation Commerciale} — Vous n'êtes pas autorisé à faire un usage commercial de cette Oeuvre, tout ou partie du matériel la composant. 
%		\item \textbf{Pas de restrictions complémentaires} — Vous n'êtes pas autorisé à appliquer des conditions légales ou des mesures techniques qui restreindraient légalement autrui à utiliser l'Oeuvre dans les conditions décrites par la licence.
%	\end{itemize}

	\end{frame}
}

\settowidth{\leftmargini}{\usebeamertemplate{itemize item}}
\addtolength{\leftmargini}{\labelsep}

\AtBeginDocument{
	\newcolumntype{L}[2]{>{\vbox to #2\bgroup\vfill\flushleft}p{#1}<{\egroup}} 
	
	\begin{frame}[plain]
		\maketitle
	\end{frame}

	\insertlicence
}


% needs etoolbox; to break links after -
\appto\UrlBreaks{\do\-}


\makeatletter
\newcommand{\xRightarrow}[2][]{\ext@arrow 0359\Rightarrowfill@{#1}{#2}}
\makeatother


\usefonttheme{structurebold}
%\usefonttheme{professionalfonts}


% figure caption
\setlength\abovecaptionskip{2pt}

% equations
\AtBeginDocument{
\abovedisplayskip=0pt
\abovedisplayshortskip=0pt
\belowdisplayskip=0pt
\belowdisplayshortskip=0pt
}


\subtitle[ML for NLP]{Chapter 02\\Machine learning for NLP} 

\changegraphpath{../img/ml4nlp/}

\begin{document}
	

\begin{frame}
	\frametitle{\inserttitle}
	\framesubtitle{ML: Example (Simple Perceptron for binary logical functions)}
	
	\begin{minipage}{0.30\textwidth} 
		\hgraphpage{perceptron.pdf}
	\end{minipage}
	%
	\begin{minipage}{0.59\textwidth}
		\scriptsize
		\begin{itemize}
			\item $ z^{(i)} = x_1^{(i)} w_1 + x_2^{(i)} w_2 + b $
			\item $ \hat{y}^{(i)} = \phi(z^{(i)}) = \begin{cases}
				1 & \text{if } z^{(i)} \ge 0 \\
				0 & \text{otherwise}
			\end{cases} $
			\item $ w_j = w_j + \nabla w_j$ where $ \nabla w_j = \frac{1}{M}\sum_{i=0}^{M} (y^{(i)} - \hat{y}^{(i)}) * x_j^{(i)} $
		\end{itemize}
	\end{minipage}
	
	\vfill
	
	\begin{exampleblock}{Training a Perceptron for AND function}
		\begin{minipage}{0.2\textwidth} 
			\scriptsize
			\begin{tabular}{|c|c|c|}
				\hline
				x\textsubscript{1} & x\textsubscript{2} & y \\
				\hline
				0 & 0 & 0  \\
				\hline
				0 & 1 & 0 \\
				\hline
				1 & 0 & 0 \\
				\hline
				1 & 1 & 1 \\
				\hline
			\end{tabular}
		\end{minipage}
		%
		\begin{minipage}{0.79\textwidth}
			\scriptsize
			\begin{itemize}
				\item Initially, $ W = [w_1, w_2, b] = [0, 0, 0] $.
				
				\item It1: $ Z = \hat{Y} = [0, 0, 0, 0] $, 
				$ \nabla W = [\frac{1}{4}, \frac{1}{4}, \frac{1}{4}] $, 
				$ W = [\frac{1}{4}, \frac{1}{4}, \frac{1}{4}] $.
				
				\item It2: $ Z = [\frac{1}{4}, \frac{2}{4}, \frac{2}{4}, \frac{3}{4}] $, 
				$ \hat{Y} = [1, 1, 1, 1] $, 
				$ \nabla W = [\frac{-1}{4}, \frac{-1}{4}, \frac{-3}{4}] $,
				$ W = [0, 0, \frac{-1}{2}] $.
				
				\item It3: $ Z = [\frac{-1}{2}, \frac{-1}{2}, \frac{-1}{2}, \frac{-1}{2}] $, 
				$ \hat{Y} = [1, 1, 1, 1] $, 
				$ \nabla W = [\frac{1}{4}, \frac{1}{4}, \frac{1}{4}] $,
				$ W = [\frac{1}{4}, \frac{1}{4}, \frac{-1}{4}] $.
				
				\item It4: $ Z = [\frac{-1}{4}, 0, 0, \frac{1}{4}] $, 
				$ \hat{Y} = Y = [0, 0, 0, 1] $ (STOP)
			\end{itemize}
		\end{minipage}
	\end{exampleblock}
	
\end{frame}
	
\begin{frame}
	\frametitle{\inserttitle}
	\framesubtitle{ML: Algorithms and abbreviations}
	
	\begin{itemize}
		\item Traditional ML algorithms
		\begin{itemize}
			\item \optword{NB}: Naive Bayes
			\item \optword{LR}: Logistic regression
			\item \optword{SVM}: Support-Vector Machine
			\item \optword{HMM}: Hidden Markov Model
			\item \optword{MEMM}: Maximum-Entropy Markov Model
			\item \optword{CRF}: Conditional Random Field
		\end{itemize}
		\item \optword{FFNN}: Feed Forward Neural Network
		\begin{itemize}
			\item \optword{MLP}: Multi-Layer Perceptron
			\item \optword{CNN}: Convolutional Neural Network
		\end{itemize} 
		\item \optword{RNN}: Recurrent Neural Network
	\end{itemize}
	
\end{frame}

\begin{frame}
	\frametitle{\inserttitle}
	\framesubtitle{ML: Types of models}
	
	\scriptsize
	\begin{tblr}{
			colspec = {p{.12\textwidth}lp{.37\textwidth}lp{.38\textwidth}},
			row{odd} = {lightblue},
			row{1} = {darkblue, fg=white, font=\bfseries, valign=m, halign=c},
			column{2,4}={white},
			column{1}={bg=darkblue, fg=white, font=\bfseries, valign=m},
			row{even} = {white},
			cell{1}{1}={white},
			colsep=3pt,
			rowsep=3pt,
			stretch = 0,
		}
		
		&& Generative && Discriminative \\
		
		&&&&\\
		
		Text && NB && LR, SVM, MLP, RNN, CNN \\
		
		&&&&\\
		
		Sequence && HMM  && MEMM, CRF, RNN \\
		
	\end{tblr}
	
	\vfill
	
	\begin{itemize}
		\item \optword{Generative model}: a model that learns to generate features given a class:
		\[\hat{Y} = \arg\max_k P(Y_k) P(X | Y_k)\]
		
		\item \optword{Discriminative model}: a model that learns to estimate a class given some features: 
		\[\hat{Y} = \arg\max_k P(Y_k | X)\]
	\end{itemize}
	
\end{frame}


\begin{frame}
	\frametitle{\inserttitle}
	\framesubtitle{NLP: Text classification}
	
	\begin{itemize}
		\item Entire text classifying:
		\begin{itemize}
			\item \optword{Sentiment analysis}: text (such as tweets) classifying into positive, negative or neutral sentiment.
			\item \optword{Spam detection}: message classifying as spam ou not spam.
			\item \optword{Figurative language detection}: text classifying as metaphor, irony, etc.
		\end{itemize}
		\item One sentence at the time classifying:
		\begin{itemize}
			\item \optword{Automatic text summarization}: sentence classifying as belonging to the summary or not.
			\item \optword{Textual implication}: checking whether a sentence implies another.
		\end{itemize}
	\end{itemize}
	
\end{frame}


\begin{frame}
	\frametitle{\inserttitle}
	\framesubtitle{NLP: Sequence classification}
	
	\begin{itemize}
		\item Given a sentence, each word must be classified separately.
		\item Example, \expword{Assign a part of speech to each word}.
		\item The word, as well as its class, can depend on the preceding (or/and following) words.
		\item \textbf{Problem}: how to classify a sequence of words that have the same (non-atomic) class?
		\item Example, \expword{[\textit{Abdelkrime Aries}]\textsubscript{\bfseries PERSON} is a teacher at the [\textit{Ecole nationale Supérieure d'Informatique}]\textsubscript{\bfseries INSTITUTE}}
	\end{itemize}
	
\end{frame}


\begin{frame}
	\frametitle{\insertshortsubtitle: \insertsection}
	\framesubtitle{NLP: Sequence classification (IOB representation)}
	
	\begin{itemize}
		\item \keyword{IOB}: Inside, Outside, Beginning
		\item For each class \textbf{CLS}, two different classes (\textbf{B-CLS} and \textbf{I-CLS}) are created.
		\item \optword{B-CLS} marks the beginning of the class.
		\item \optword{I-CLS} marks the continuation of the class.
		\item \optword{O} means "no class".
		\item Example, \expword{Abdelkrime/\textsubscript{\bfseries B-PERSON} Aries/\textsubscript{\bfseries I-PERSON}  is/\textsubscript{\bfseries O} a/\textsubscript{\bfseries O} teacher/\textsubscript{\bfseries O} at/\textsubscript{\bfseries O} the/\textsubscript{\bfseries O} Ecole/\textsubscript{\bfseries B-INSTITUTE} nationale/\textsubscript{\bfseries I-INSTITUTE} Supérieure/\textsubscript{\bfseries I-INSTITUTE} d'/\textsubscript{\bfseries I-INSTITUTE} Informatique/\textsubscript{\bfseries I-INSTITUTE}}.
	\end{itemize}
\end{frame}


\begin{frame}
	\frametitle{\inserttitle}
	\framesubtitle{NLP: Text generation}
	
%	TODO 
	\begin{itemize}
		\item 
	\end{itemize}
	
\end{frame}


\begin{frame}
	\frametitle{\inserttitle}
	\framesubtitle{\insertshortsubtitle: Plan}
	
	\begin{multicols}{2}
		\small
		\tableofcontents
	\end{multicols}
\end{frame}



%===================================================================================
\section{Traditional ML and MLP-only architectures}
%===================================================================================

\begin{frame}
	\frametitle{\insertshortsubtitle}
	\framesubtitle{\insertsection}
	
	\begin{itemize}
		\item 
	\end{itemize}
	
\end{frame}


\subsection{Text classification}

\begin{frame}
	\frametitle{\insertshortsubtitle: \insertsection}
	\framesubtitle{\insertsubsection: Examples}
	
	\vgraphpage{text_class_trad_exp.pdf}
	
\end{frame}

\subsection{HMM for Sequence classification}


\begin{frame}
	\frametitle{\insertshortsubtitle: \insertsection}
	\framesubtitle{\insertsubsection: Markov model}
	
	\begin{minipage}{.64\textwidth}
			\begin{itemize}
					\item Markov hypothesis
					\item $ P(q_i = a | q_1, \ldots, q_{i-1}) \approx P(q_i = a | q_{i-1}) $
					\item $Q = \{q_1, q_2, \ldots, q_n\}$: stats.
					\item $A = \begin{bmatrix}%
							a_{11} & a_{12} & \ldots & a_{1n} \\
							a_{21} & a_{22} & \ldots & a_{2n} \\
							\vdots & \vdots & \ddots & \vdots \\
							a_{n1} & a_{n2} & \ldots & a_{nn} \\
						\end{bmatrix}$: transition probability matrix.
					\item $\sum_j a_{ij} = 1,\, \forall i$.
				\end{itemize}
		\end{minipage}
	\begin{minipage}{.35\textwidth}
			\hspace*{-1cm}
			\begin{tikzpicture}[
					> = stealth, % arrow head style
					shorten > = 1pt, % don't touch arrow head to node
					auto,
					node distance = 2cm, % distance between nodes
					semithick, % line style
					font=\tiny\bfseries
					]
					
					\node[circle,draw] (qC) {C};
					\node[circle,draw] (qH) [below left of=qC] {H};
					\node[circle,draw] (qW) [below right of=qC] {W};
					
					\path[->] 	
					(qC) 	edge [loop above] node {0.8} ()
					edge [] node {0.1} (qH)
					edge [bend left] node {0.1} (qW)
					(qH) 	edge [loop left] node {0.6} ()
					edge [bend left] node {0.1} (qC)
					edge [] node {0.3} (qW)
					(qW)	edge [loop right] node {0.6} ()
					edge [bend left] node {0.3} (qH)
					edge [] node {0.1} (qC);
				\end{tikzpicture}
		\end{minipage}
	
	\begin{itemize}
			\item $\pi = [\pi_1, \pi_2, \ldots, \pi_n ]$: states initial probability distribution.
			\item $\sum_i \pi_i = 1$.
			\vspace{-6pt}\item E.g. \expword{Calculate $P(H\, W\, C\, C)$ where $\pi = [\overbrace{0.1}^C, \overbrace{0.7}^H, \overbrace{0.2}^W]$}.
		\end{itemize}
\end{frame}

\begin{frame}
	\frametitle{\insertshortsubtitle: \insertsection}
	\framesubtitle{\insertsubsection: Hidden Markov model}
	\begin{minipage}{.54\textwidth}
			\begin{itemize}
					\item $Q = \{q_1, q_2, \ldots, q_n\}$: states.
					\item $A$: transition probability matrix where $\sum_j a_{ij} = 1,\, \forall i$.
					\item $O = o_1 o_2 \ldots o_T$: sequence of observed events (words).
					\item $B = b_i(o_t)$: observation probabilities (\keyword{emission probabilities}), each represents the probability of generating an observation $o_t$ in a state $q_i$.
				\end{itemize}
		\end{minipage}
	\begin{minipage}{.45\textwidth}
			\begin{tikzpicture}[
					> = stealth, % arrow head style
					shorten > = 1pt, % don't touch arrow head to node
					auto,
					node distance = 1.5cm, % distance between nodes
					semithick, % line style
					font=\bfseries\fontsize{4}{4}\selectfont
					]
					
					\node[circle,draw] (q2) {MD};
					\node[align=center,draw] (q2e) [left of=q2] {B2\\ \\P(``will"|MD)\\...\\P(``the"|MD)\\...\\P(``back"|MD)};
					\node[circle,draw] (q3) [right of=q2] {NN};
					\node[align=center,draw] (q3e) [right of=q3] {B3\\ \\P(``will"|NN)\\...\\P(``the"|NN)\\...\\P(``back"|NN)};
					\node[circle,draw] (q1) [below of=q2] {VB};
					\node[align=center,draw] (q1e) [left of=q1] {B1\\ \\P(``will"|VB)\\...\\P(``the"|VB)\\...\\P(``back"|VB)};
					
					\path[->] 	
					(q1) 	edge [loop below] node {a11} ()
					edge [bend left] node {a12} (q2)
					edge [] node {a13} (q3)
					(q2) 	edge [loop above] node {a22} ()
					edge [bend left] node {a21} (q1)
					edge [bend right] node {a23} (q3)
					(q3)	edge [loop above] node {a33} ()
					edge [bend left] node {a31} (q1)
					edge [bend right] node {a32} (q2);
					
					\path[dashed] 	
					(q1) 	edge [] node {} (q1e)
					(q2) 	edge [] node {} (q2e)
					(q3) 	edge [] node {} (q3e);
					
				\end{tikzpicture}
		\end{minipage}
	
	\begin{itemize}
			\item $\pi = [\pi_1, \pi_2, \ldots, \pi_n ]$: states initial probability distribution.
			\item $\sum_i \pi_i = 1$.
		\end{itemize}
	
\end{frame}

\begin{frame}
	\frametitle{\insertshortsubtitle: \insertsection}
	\framesubtitle{\insertsubsection: Training}
	
	\begin{itemize}
			\item States are tags (categories) $t_i$.
			\item Let's define \keyword{START} as sentences' start tag.
			\item The observations $o_i$ are the words $w_i$.
			\item \keyword{C}: Number of occurrences in the training corpus.
		\end{itemize}
	
	\begin{block}{HMM training}
			\[
			\text{Transition probabilities: } P(t_i | t_{i-1}) = \frac{C(t_{i-1}, t_i)}{C(t_{i-1})} 
			\]\[
			\text{Emission probabilities: } P(w_i | t_i) = \frac{C(t_i, w_i)}{C(t_i)}
			\]\[
			\text{Initial distribution: } \pi_i = P(t_i | START) = \frac{C(START, t_i)}{C(START)}
			\]
		\end{block}
\end{frame}

\begin{frame}
	\frametitle{\insertshortsubtitle: \insertsection}
	\framesubtitle{\insertsubsection: Labeling (tag decoding)}
	
	\begin{itemize}
			\item Given
			\begin{itemize}
					\item a Markov model $\lambda = (A, B)$,
					\item an observations' sequence (words): $w = w_1 w_2 \ldots w_n$
				\end{itemize}
			\item Estimate a labels' sequence $\hat{t} = \hat{t}_1 \hat{t}_2 \ldots \hat{t}_n$
		\end{itemize}
	
	\begin{block}{Labels' decoding using HMM}
			\[
			\hat{t} = \arg\max\limits_t P(t | w) = \arg\max\limits_t \frac{P(w|t) P(t)}{P(w)} = \arg\max\limits_t P(w|t) P(t)%\text{ tel que } t = t_1 t_2 \ldots t_n
			\]
			
			\[ 
			P(w | t) \approx \prod\limits_{i=1}^n P(w_i|t_i) 
			\hskip2cm
			P(t) \approx \prod\limits_{i=1}^n P(t_i|t_{i-1}) 
			\]
			
			\[
			\hat{t} = \arg\max\limits_t \prod\limits_{i=1}^n P(w_i|t_i) P(t_i|t_{i-1})
			\]
		\end{block}
\end{frame}

\begin{frame}
	\frametitle{\insertshortsubtitle: \insertsection}
	\framesubtitle{\insertsubsection: Viterbi algorithm}
	
	\begin{block}{Viterbi}
			\scriptsize
			\begin{algorithm}[H]
		%				\vskip-1em
					\KwData{$w = w_1 \ldots w_T$, HMM $\lambda = (A, B)$ with $N$ states}
					\KwResult{$best\_path$, $prob\_path$}
					
					Create a matrix $viterbi[N, T]$\;
					
					\lForEach{state $ s = 1 \ldots N$}{
							$viterbi[s, 1] = \pi_s * b_s(w_1);\, backpointer[s, 1] = 0$
						}
					
					\ForEach{tag $ t = 2 \ldots T$}{
							\ForEach{state $ s = 1 \ldots N$}{
									$viterbi[s, t] = \max\limits_{s'=1}^N viterbi[s', t-1] * a_{s',s} * b_s(w_t)$\;
									$backpointer[s, t] = \arg\max\limits_{s'=1}^N viterbi[s', t-1] * a_{s',s} * b_s(w_t)$\;
								}
						}
					
					$prob\_path = \max\limits_{s=1}^N viterbi[s, T];\, pointer\_path = \arg\max\limits_{s=1}^N viterbi[s, T]$\;
					
					$best\_path$ is the path starting from $pointer\_path$ and following $backpointer$
					
					\Return $best\_path$, $prob\_path$\;
					
				\end{algorithm}
		\end{block}
	
\end{frame}

\subsection{MEMM for for Sequence classification}

\begin{frame}
	\frametitle{\insertshortsubtitle: \insertsection}
	\framesubtitle{\insertsubsection}
	
	\begin{itemize}
			\item Let $x_{n}^{m} = x_n x_{n+1} \ldots x_{m-1} x_m$
			\item Given 
			\begin{itemize}
					\item an observations sequence (words): $w = w_1 w_2 \ldots w_n$
					\item a set of features $f$ defined on the sequences (such as: majuscule word)
				\end{itemize}
			\item Estimate tags sequence $\hat{t} = \hat{t}_1 \hat{t}_2 \ldots \hat{t}_n$
		\end{itemize}
	
	\begin{block}{Labels decoding using MEMM}
			\[
			\hat{t} = \arg\max\limits_t P(t | w) = \arg\max\limits_t \prod\limits_{i}  P(t_i | w_i, t_{i-1})
			\]
			
			\[
			\hat{t} = \arg\max\limits_t \prod\limits_{i}  
			\frac{exp\left(\sum_j \theta_j f_j(t_i, w_{i-l}^{i+l}, t_{i-k}^{i-1})\right)}%
			{\sum_{t' \in tags} exp\left(\sum_j \theta_j f_j(t'_i, w_{i-l}^{i+l}, t_{i-k}^{i-1})\right)}
			\]
		\end{block}
	
\end{frame}

% TODO
\subsection{Text generation}

\begin{frame}
	\frametitle{\insertshortsubtitle: \insertsection}
	\framesubtitle{\insertsubsection}
	
\end{frame}


%===================================================================================
\section{CNN-based architectures}
%===================================================================================

\begin{frame}
	\frametitle{\insertshortsubtitle}
	\framesubtitle{\insertsection}
	
	\begin{itemize}
		\item 
	\end{itemize}
	
\end{frame}

\subsection{Text classification}

\begin{frame}
	\frametitle{\insertshortsubtitle: \insertsection}
	\framesubtitle{\insertsubsection}
	
	\begin{itemize}
			\item For each word of the text, a vector representation must be defined (vector of \keyword{N} elements).
			\item Representations are concatenated vertically (\keyword{M} words).
			\item This will result in a matrix of \keyword{M X N} elements.
			\item A one-dimensional convolution (dimension = \keyword{K}) is applied; i.e. a \keyword{K X N} dimension filter.
			\item Several filters can be used to learn several representations.
			\item Pooling can also be used.
			\item A \keyword{MLP} layer (or several) is used at the end to estimate the class.
			\item \textbf{Problem}: Technically, we cannot define a neural network with a variable number of words.
			\item \textbf{Solution}: \textcolor{yellow!30}{A maximum number must be defined; if the number of words is lower, padding is added; otherwise, extra words are deleted.}
		\end{itemize}
	
\end{frame}

\begin{frame}
	\frametitle{\insertshortsubtitle: \insertsection}
	\framesubtitle{\insertsubsection: Examples}
	
	\vgraphpage{text_class_CNN_exp.pdf}
	
\end{frame}

%TODO CNN for Sequence classification
\subsection{Sequence classification}


% TODO CNN for Text generation
\subsection{Text generation}

\begin{frame}
	\frametitle{\insertshortsubtitle: \insertsection}
	\framesubtitle{\insertsubsection}
	
\end{frame}


%===================================================================================
\section{RNN-based architectures}
%===================================================================================

\begin{frame}
	\frametitle{\insertshortsubtitle}
	\framesubtitle{\insertsection}
	
	\begin{itemize}
		\item 
	\end{itemize}
	
\end{frame}

\subsection{Text classification}

\begin{frame}
	\frametitle{\insertshortsubtitle: \insertsection}
	\framesubtitle{\insertsubsection}
	
	\begin{itemize}
			\item For each word of the text, a vector representation must be defined (vector of \keyword{N} elements).
			\item The RNN network creates M cells corresponding to the \keyword{M} words.
			\item Each cell calculates an intermediate state (vector) and passes it to the next cell.
			\item The last cell generates a representation of the text (the probability of the last word's occurrence given the past words).
			\item This representation is introduced to a \keyword{MLP} to estimate the class of the text.
			\item \textbf{Problem}: Technically, during training, words of different sizes cannot be processed with a single matrix operation.
			\item \textbf{Solution}: \textcolor{yellow!30}{A maximum number must be defined; if the words number is lower, padding is added; otherwise, extra words are deleted}
		\end{itemize}
	
\end{frame}

\begin{frame}
	\frametitle{\insertshortsubtitle: \insertsection}
	\framesubtitle{\insertsubsection: Examples}
	
	\vgraphpage{text_class_RNN_exp.pdf}
	
\end{frame}


\subsection{Sequence classification}

\begin{frame}
	\frametitle{\insertshortsubtitle: \insertsection}
	\framesubtitle{\insertsubsection: Example}
	
	\vgraphpage{seq_class_RNN_exp.pdf}
	
\end{frame}

% TODO RNN for Text generation
\subsection{Text generation}

\begin{frame}
	\frametitle{\insertshortsubtitle: \insertsection}
	\framesubtitle{\insertsubsection}
	
	Language model
	GAN-based decoder
	
	
\end{frame}


%===================================================================================
\section{Transformer-based architectures}
%===================================================================================

% TODO
\begin{frame}
	\frametitle{\insertshortsubtitle}
	\framesubtitle{\insertsection}
	
	\begin{itemize}
		\item 
	\end{itemize}
	
\end{frame}

% TODO Transformers for Text classification
\subsection{Text classification}

\begin{frame}
	\frametitle{\insertshortsubtitle: \insertsection}
	\framesubtitle{\insertsubsection}
	
	encoder-based
	decoder-based
	full-transformer
	
\end{frame}

% TODO Transformers for Sequence classification
\subsection{Sequence classification}

\begin{frame}
	\frametitle{\insertshortsubtitle: \insertsection}
	\framesubtitle{\insertsubsection}
	
	encoder-based
	decoder-based
	full-transformer
	
\end{frame}


% TODO Transformers for Text generation
\subsection{Text generation}

\begin{frame}
	\frametitle{\insertshortsubtitle: \insertsection}
	\framesubtitle{\insertsubsection}
	
	encoder-based
	decoder-based
	full-transformer
	
\end{frame}



\insertbibliography{NLP02}{*}
	
\end{document}

